\section{Related Work}
\label{sec:related}

\textsc{ui-bench} sits at the intersection of three established lines of work:
visual program induction, synthetic diagnostic benchmarks, and image-to-code
evaluation of multimodal models.

\paragraph{Visual program induction and inverse graphics.}
Predicting executable programs from images has a long history. CSGNet~\cite{sharma2018csgnet}
infers constructive solid geometry programs from 2D and 3D shapes, and is the
direct conceptual ancestor of \textsc{ui-bench}. \cite{liu2019scenes} learn to
describe scenes with a DSL that supports loops and grouping, demonstrating that
compositional program structure -- and not only local shape identity -- can be
recovered from a single image. \cite{li2020perspective} and
\cite{li2020multiplane} extend program induction to perspective scenes and
repeated 3D structure, while \cite{duan2022parametric} studies parametric
primitives with explicit function correlations, close in spirit to the
integer-parameter DSL of \textsc{ui-bench}. LILO~\cite{grand2024lilo} synthesizes
reusable program libraries across domains including graphics composition,
providing a template for symbolic baselines. These systems prove that the
image-to-program problem is well-defined; none of them are benchmark-first
evaluations of modern multimodal models.

\paragraph{Benchmark design.}
\textsc{CLEVR}~\cite{johnson2017clevr} showed how much scientific value a
synthetic, carefully factorized benchmark can add when it is designed to reduce
spurious shortcuts and expose reasoning failure modes.
\textsc{CLOSURE}~\cite{bahdanau2019closure} extended this insight by probing
systematic generalization. We adopt the same diagnostic posture: rather than
building a large, noisy dataset, we expose the axes of variation explicitly and
allow researchers to regenerate the dataset from a seed.

\paragraph{Closest benchmark predecessors.}
\textsc{TurtleBench}~\cite{rismanchian2025turtlebench} evaluates vision-language
models on turtle-geometry programs and is the closest benchmark-level neighbor.
It reports that strong models still struggle, which supports the general thesis
that visual program reconstruction is hard. \textsc{ui-bench} differs along three
axes: (1) the DSL is a tiny shape-primitive set rather than turtle paths, which
removes path-planning reasoning and isolates perception-plus-emission; (2) the
benchmark is explicitly renewable via fresh seeds; and (3) scoring is raster-based
and deterministic.

\textsc{Image2Struct}~\cite{roberts2024image2struct} introduced round-trip
structure extraction: image $\rightarrow$ structure $\rightarrow$ rendered image
$\rightarrow$ similarity. Our scoring pipeline follows the same philosophy.
Unlike \textsc{Image2Struct}, which spans noisy real-world web pages, LaTeX, and
music, \textsc{ui-bench} stays inside a small controlled space so that reported
failures can be cleanly attributed to perception or emission.

\paragraph{Broader image-to-code.}
pix2code~\cite{beltramelli2017pix2code}, Design2Code~\cite{si2024design2code},
\textsc{VCode}~\cite{lin2025vcode}, and \textsc{Omni-I2C}~\cite{zhou2026omni}
demonstrate that image-to-code is now a mature benchmark family spanning GUI
screenshots, SVG, and general graphics-to-code. These benchmarks mix many
confounders: OCR, library conventions, rendering-engine variability, and external
assets. \textsc{ui-bench} strips away these confounders by design, providing a
sibling benchmark whose strengths are control and reproducibility rather than
realism.

\paragraph{How to read the contribution.}
\textsc{ui-bench} is not the first benchmark to ask models to reconstruct
programs from images. Its contribution is a \emph{specific combination}:
a deterministic render-based evaluator, explicit difficulty axes, a
provider-agnostic adapter for cost-controlled runs, and a renewable frozen
evaluation split with publicly verifiable instance hashes. We view it as
complementary to \textsc{TurtleBench} and \textsc{Image2Struct} rather than a
replacement.
